\section{Conclusion and Discussion}

This paper studied a lazy MUS-guided greedy hitting-set algorithm for the
Minimum-Weight Minimum Correction Subset (MWMCS) problem on synthetic
HM-style type-equality constraint systems. Across 300 instances spanning three
generator families and five sizes, the greedy algorithm matched the exact Z3
optimum on every instance, with approximation ratio uniformly equal to one.
The brute-force baseline independently confirmed the Z3 optima on small
instances. Runtime analysis showed that lazy greedy overhead is driven by the
number of MUSes discovered rather than by the total constraint count: Datasets
A and C, each with one MUS per instance, ran at or below exact runtime, while
Dataset~B's five-iteration loop produced a consistent $2$--$3\times$ slowdown
that was stable across instance sizes.

\paragraph{Interpretation.}
The absence of an approximation gap describes the tested instance space, not
the algorithm. A greedy hitting-set algorithm carries a worst-case bound of
$1 + \ln d$, where $d$ is the maximum MUS size, and this bound remains
applicable to instances outside the tested space. The empirical result here
is that the three synthetic generators produce structurally simple MUS
hypergraphs---a single small MUS in Datasets A and C, and a regular clique
on a single shared variable in Dataset~B---on which weighted greedy happens
to be optimal. The experiment therefore characterises the generators as much
as it characterises the algorithm.

\paragraph{Limitations.}
Three limitations should be stated explicitly. First, the synthetic
generators are HM-shaped rather than HM-derived: they encode plausible
type-equality structure but were not produced by a real Haskell or OCaml
compiler. Second, only one greedy variant was evaluated; alternative
strategies, such as weighted greedy with full MUS enumeration or LP-rounding,
were not compared. Third, the instance sizes ($|C| \le 200$) are small
relative to real-world programs, where MUS hypergraph structure is likely
richer.

\paragraph{Future work.}
The most direct extension is to construct adversarial generators that
exercise the worst-case greedy behaviour, for example by planting multiple
disjoint MUSes with cross-cutting weight gradients in the style of the Set
Cover hardness template of Feige~\cite{feige1998}. A complementary direction,
returning to the theoretical question framed in our preliminary
proposal phase, is to establish whether logarithmic inapproximability
transfers from Set Cover to MWMCS on HM-derived instance classes, or whether
the AST-derived structure of HM constraints admits strictly better bounds.
Empirical evaluation on real type-error corpora, contingent on access to a
reusable benchmark, remains an open opportunity.